% ============================================================
%  NOTATION TABLE
%  Placed after the Table of Contents in frontmatter.
% ============================================================

\chapter*{Обозначения}
\addcontentsline{toc}{chapter}{Обозначения}

В следующих таблицах собраны символы, используемые на протяжении всей книги единообразно. Если символ несёт значение, специфическое для отдельной главы, эта глава вводит его явно; приведённые здесь записи охватывают глобально повторяющиеся соглашения.

% ----------------------------------------------------------------
\section*{Общая математика}
% ----------------------------------------------------------------

\begin{center}
\begin{tabular}{@{}ll@{}}
\toprule
\textbf{Символ} & \textbf{Значение} \\
\midrule
$\R$, $\R^n$ & Вещественные числа; $n$-мерное вещественное векторное пространство \\
$\N$         & Натуральные числа $\{0,1,2,\ldots\}$ \\
$\Z$         & Целые числа \\
$[n]$        & Индексное множество $\{1,2,\ldots,n\}$ \\
$|S|$        & Мощность множества $S$ \\
$\langle x, y\rangle$ & Скалярное произведение векторов $x$ и $y$ \\
$\|x\|$, $\|x\|_p$   & Евклидова норма (2-норма); $\ell^p$-норма \\
$\|A\|_F$    & Норма Фробениуса матрицы $A$ \\
$A^\top$     & Транспонирование матрицы $A$ \\
$A^{-1}$     & Обратная матрица к квадратной матрице $A$ \\
$\nabla_\theta f$ & Градиент $f$ по $\theta$ \\
$\nabla^2_\theta f$ & Гессиан $f$ по $\theta$ \\
$\partial_x f$ & Частная производная $f$ по $x$ \\
$f \circ g$  & Композиция функций: $(f\circ g)(x)=f(g(x))$ \\
$\propto$    & Пропорционально (с точностью до нормировочной константы) \\
$\ind\{A\}$  & Индикаторная функция: $1$, если событие $A$ произошло, иначе $0$ \\
$\log$       & Натуральный логарифм (основание $e$) \\
$\log_2$     & Двоичный логарифм (основание $2$) \\
$\exp(x)$    & Показательная функция $e^x$ \\
$[x]_+$      & Положительная часть: $\max(x,0)$ \\
$\lfloor x \rfloor$, $\lceil x \rceil$ & Функции пола и потолка \\
\bottomrule
\end{tabular}
\end{center}

% ----------------------------------------------------------------
\section*{Теория вероятностей и статистика}
% ----------------------------------------------------------------

\begin{center}
\begin{tabular}{@{}ll@{}}
\toprule
\textbf{Символ} & \textbf{Значение} \\
\midrule
$\Prob(A)$          & Вероятность события $A$ \\
$\Prob(A\mid B)$    & Условная вероятность $A$ при условии $B$ \\
$\E[X]$             & Математическое ожидание случайной величины $X$ \\
$\E[X \mid Y]$      & Условное математическое ожидание $X$ при условии $Y$ \\
$\Var[X]$           & Дисперсия $X$: $\E[X^2]-(\E[X])^2$ \\
$\Cov[X,Y]$         & Ковариация $X$ и $Y$ \\
$p(x)$              & Функция плотности вероятности или функция масс вероятности \\
$p(x \mid y)$       & Условная плотность/масса $x$ при условии $y$ \\
$X \sim p$          & $X$ распределена по закону $p$ \\
$X \perp Y$         & $X$ и $Y$ независимы \\
$\cN(\mu,\sigma^2)$ & Нормальное распределение со средним $\mu$ и дисперсией $\sigma^2$ \\
$\mathrm{Bernoulli}(p)$ & Распределение Бернулли с вероятностью успеха $p$ \\
$\mathrm{Cat}(\mathbf{p})$ & Категориальное распределение с вектором вероятностей $\mathbf{p}$ \\
$\KL{p}{q}$         & Дивергенция Кульбака-Лейблера от $q$ к $p$ \\
$H(p)$              & Энтропия Шеннона распределения $p$ \\
$H(p,q)$            & Перекрёстная энтропия между $p$ и $q$ \\
$I(X;Y)$            & Взаимная информация между $X$ и $Y$ \\
$\sigma(x)$         & Сигмоидная функция $1/(1+e^{-x})$ \\
$\mathrm{softmax}(\mathbf{z})_i$ & $e^{z_i}/\sum_j e^{z_j}$ \\
\bottomrule
\end{tabular}
\end{center}

% ----------------------------------------------------------------
\section*{Основные обозначения обучения с подкреплением}
% ----------------------------------------------------------------

\begin{center}
\begin{tabular}{@{}ll@{}}
\toprule
\textbf{Символ} & \textbf{Значение} \\
\midrule
$\cS$               & Пространство состояний \\
$\cA$               & Пространство действий \\
$\cR$               & Пространство вознаграждений (также используется для сожаления) \\
$\cM$               & Кортеж МДП $(\ cS,\cA,P,R,\gamma)$ \\
$P(s'\mid s,a)$     & Функция вероятности переходов \\
$R(s,a)$, $r(s,a)$  & Ожидаемое вознаграждение; сигнал вознаграждения \\
$s_t, a_t, r_t$     & Состояние, действие, вознаграждение на шаге $t$ \\
$\tau$              & Траектория (эпизод): $(s_0,a_0,r_0,s_1,\ldots)$ \\
$T$                 & Горизонт (длина эпизода); может быть $\infty$ \\
$\gamma \in [0,1)$  & Коэффициент дисконтирования \\
$G_t$               & Доход с момента $t$: $\sum_{k=0}^{\infty}\gamma^k r_{t+k+1}$ \\
$\pi(a\mid s)$      & Стохастическая стратегия \\
$\pi(s)$            & Детерминированная стратегия (отображает $s$ в действие) \\
$\pi^*$             & Оптимальная стратегия \\
$V^\pi(s)$          & Функция ценности состояния при стратегии $\pi$ \\
$Q^\pi(s,a)$        & Функция ценности действия (Q-функция) при стратегии $\pi$ \\
$V^*(s)$, $Q^*(s,a)$ & Оптимальные функции ценности \\
$A^\pi(s,a)$        & Функция преимущества $Q^\pi(s,a)-V^\pi(s)$ \\
$\delta_t$          & Ошибка временной разности на шаге $t$ \\
$\alpha$            & Скорость обучения / размер шага \\
$\epsilon$          & Коэффициент исследования (например, $\epsilon$-жадный) \\
$\lambda$           & Скорость затухания следов приемлемости (TD($\lambda$)) \\
$n$                 & $n$-шаговый просмотр вперёд \\
$\theta, \phi, \psi$  & Векторы параметров (стратегии, функции ценности, модели вознаграждения) \\
$\pi_\theta$        & Стратегия, параметризованная $\theta$ \\
$V_\phi$, $Q_\phi$  & Функция ценности, параметризованная $\phi$ \\
$\mathcal{T}^\pi$   & Оператор Беллмана для стратегии $\pi$ \\
$\mathcal{T}^*$     & Оператор оптимальности Беллмана \\
$d^\pi(s)$          & Стационарное распределение состояний при стратегии $\pi$ \\
$J(\theta)$         & Целевая функция стратегии (ожидаемый доход) \\
$\hat{A}_t$         & Оценка преимущества на шаге $t$ \\
$\hat{g}$           & Стохастическая оценка градиента стратегии \\
$L^\mathrm{CLIP}$   & Усечённая суррогатная целевая функция PPO \\
$\rho_t(\theta)$    & Отношение вероятностей $\pi_\theta(a_t|s_t)/\pi_{\theta_\mathrm{old}}(a_t|s_t)$ \\
\bottomrule
\end{tabular}
\end{center}

% ----------------------------------------------------------------
\section*{Обозначения для языковых моделей}
% ----------------------------------------------------------------

\begin{center}
\begin{tabular}{@{}ll@{}}
\toprule
\textbf{Символ} & \textbf{Значение} \\
\midrule
$\cV$                    & Словарь (множество токенов) \\
$x$                      & Промпт / входная последовательность \\
$y$                      & Ответ / выходная последовательность \\
$y_t$                    & $t$-й токен ответа $y$ \\
$y_{<t}$                 & Токены ответа до позиции $t$ \\
$|y|$                    & Длина ответа $y$ (в токенах) \\
$\mathrm{eos}$           & Токен конца последовательности \\
$\pi_\theta(y\mid x)$    & Стратегия языковой модели (авторегрессионная) \\
$\pi_\mathrm{ref}(y\mid x)$ & Эталонная (референсная) стратегия \\
$r_\psi(x,y)$            & Модель вознаграждения с параметрами $\psi$ \\
$\beta$                  & Коэффициент KL-штрафа \\
$\cD_\mathrm{pref}$      & Набор данных о предпочтениях людей \\
$y^+$, $y^-$             & Предпочтительный и отвергнутый ответы \\
$p^*(y^+\succ y^-)$      & Вероятность предпочтения людей \\
$Z(x)$                   & Нормировочная функция (статсумма) для оптимальной стратегии LM \\
$\ell_\mathrm{DPO}$      & Функция потерь DPO \\
$g_i$                    & $i$-е сгенерированное завершение (GRPO) \\
$\hat{r}_i$              & Нормализованное вознаграждение для группы $i$ (GRPO) \\
\bottomrule
\end{tabular}
\end{center}

% ----------------------------------------------------------------
\section*{Распространённые сокращения}
% ----------------------------------------------------------------

\begin{center}
\begin{tabular}{@{}ll@{}}
\toprule
\textbf{Сокращение} & \textbf{Расшифровка} \\
\midrule
RL     & Reinforcement Learning (обучение с подкреплением) \\
MDP    & Markov Decision Process (марковский процесс принятия решений) \\
POMDP  & Partially Observable MDP (частично наблюдаемый МДП) \\
TD     & Temporal Difference (временная разность) \\
MC     & Monte Carlo (метод Монте-Карло) \\
DP     & Dynamic Programming (динамическое программирование) \\
GPI    & Generalised Policy Iteration (обобщённая итерация стратегии) \\
VFA    & Value Function Approximation (аппроксимация функции ценности) \\
DQN    & Deep Q-Network \\
DDPG   & Deep Deterministic Policy Gradient \\
A2C/A3C & Advantage Actor-Critic / Asynchronous A3C \\
PPO    & Proximal Policy Optimisation \\
TRPO   & Trust Region Policy Optimisation \\
GAE    & Generalised Advantage Estimation \\
SGD    & Stochastic Gradient Descent (стохастический градиентный спуск) \\
SFT    & Supervised Fine-Tuning (дообучение с учителем) \\
RLHF   & Reinforcement Learning from Human Feedback \\
DPO    & Direct Preference Optimisation \\
GRPO   & Group Relative Policy Optimisation \\
RLVR   & RL with Verifiable Rewards (RL с верифицируемыми вознаграждениями) \\
CoT    & Chain-of-Thought (цепочка рассуждений) \\
LLM    & Large Language Model (большая языковая модель) \\
KL     & Kullback-Leibler (дивергенция Кульбака-Лейблера) \\
i.i.d. & Independent and Identically Distributed (независимые одинаково распределённые) \\
w.r.t. & With Respect To (по отношению к) \\
\bottomrule
\end{tabular}
\end{center}
